% Part: computability
% Chapter: computability-theory
% Section: halting-problem

\documentclass[../../../include/open-logic-section]{subfiles}

\begin{document}

\olfileid[hi]{cmp}{thy}{hlt}
\olsection{विराम समस्या}

रचना के अनुसार सार्वत्रिक आंशिक संगणनीय फलन~$\fn{Un}(e,x)$ तभी और केवल
तभी परिभाषित है, जब~$e$ से कूटित फलन की संगणना आगत~$x$ के लिए कोई मान
उत्पन्न करती है। यह पूछना स्वाभाविक है कि क्या हम तय कर सकते हैं कि ऐसा
होता है या नहीं। वास्तव में हम ऐसा नहीं कर सकते। संगणना के ट्यूरिंग मशीन
मॉडल में इसका अर्थ है कि कोई दी हुई ट्यूरिंग मशीन किसी दी हुई आगत पर
रुकती है या नहीं, यह संगणनात्मक रूप से अनिर्णेय है। इसलिए अगला प्रमेय
``विराम समस्या की अनिर्णेयता'' कहलाता है। नीचे हम दो प्रमाण देंगे। पहला
हमारी पिछली चर्चा की कड़ी को आगे बढ़ाता है, जबकि दूसरा अधिक प्रत्यक्ष है।

\begin{thm}
\ollabel{thm:halting-problem}
मानें
\[
h(e, x)  =
\begin{cases}
1 & \text{यदि\/ $\fn{Un}(e, x)$ परिभाषित है} \\
0 & \text{अन्यथा।}
\end{cases}
\]
तब $h$ संगणनीय नहीं है।
\end{thm}

\begin{proof}
मान लें कि $h$ संगणनीय है। हम दिखाते हैं कि तब हम एक सार्वत्रिक संगणनीय
फलन परिभाषित कर सकते हैं। परिभाषित करें
\[
\fn{Un'}(e,x) =
\begin{cases}
\fn{Un}(e,x) & \text{यदि $h(e,x) = 1$} \\
0 & \text{अन्यथा।}
\end{cases}
\]
अब $\fn{Un'}(e, x)$ एक पूर्ण फलन है, और यदि~$h$ संगणनीय है तो यह भी
संगणनीय है। उदाहरणतः, आदिम पुनरावर्तन से $g$ को इस प्रकार परिभाषित कर सकते हैं:
\begin{align*}
g(0, e, x) & \simeq 0 \\
g(y+1, e, x) & \simeq \fn{Un}(e,x);
\end{align*}
तब
\[
\fn{Un'}(e,x) \simeq g(h(e,x),e,x).
\]
$\fn{Un'}(e,x)$, $\fn{Un}(e,x)$ से हर उस स्थान पर सहमत है जहाँ बाद वाला
परिभाषित है; अतः $\fn{Un'}$ उन आंशिक संगणनीय फलनों के लिए सार्वत्रिक है जो पूर्ण
हैं। किंतु यह \olref[nou]{thm:no-univ} के विरुद्ध है।
\end{proof}

\begin{proof}
मान लें कि $h(e,x)$ संगणनीय होता। फलन $g$ को इस प्रकार परिभाषित करें:
\[
g(x) =
\begin{cases}
  0                & \text{यदि $h(x,x) = 0$} \\
  \fundefined & \text{अन्यथा।}
\end{cases}
\]
फलन $g$ आंशिक संगणनीय है। उदाहरणतः, इसे $\umin{y}{h(x,x) = 0}$ के रूप
में परिभाषित किया जा सकता है। अतः किसी~$e$ के लिए
$g(x) \simeq \fn{Un}(e,
x)$ प्रत्येक~$x$ पर। क्या $g$, $e$ पर परिभाषित है? यदि हाँ,
तो~$g$ की परिभाषा से $h(e,e) = 0$ ($h$~मान~$0$ तभी ले सकता है जब वह
परिभाषित हो)। फिर~$h$ की परिभाषा से $\fn{Un}(e, e)$ अपरिभाषित है।
हमारी मान्यता $g(x) \simeq \fn{Un}(e, x)$, जो प्रत्येक~$x$ के लिए है,
से इसका अर्थ है
कि $g(e)$ अपरिभाषित है—विरोधाभास। दूसरी ओर, यदि $g(e)$ अपरिभाषित है,
तो $h(e,e) \neq 0$, और इसलिए $h(e,e) = 1$। फलतः $\fn{Un}(e, e)$
परिभाषित है। किंतु $g(x) \simeq \fn{Un}(e, x)$ होने से $g(e)$ भी
परिभाषित होगा। फिर से विरोधाभास।
\end{proof}

\begin{tagblock}{TMs}
\begin{explain}
इस तर्क को ट्यूरिंग मशीनों के पदों में भी कहा जा सकता है। मान लें कि
कोई ट्यूरिंग मशीन~$H$ है, जो किसी ट्यूरिंग मशीन~$E$ का वर्णन और कोई
आगत~$x$ लेकर तय करती है कि $E$, आगत~$x$ पर रुकती है या नहीं। तब हम एक
दूसरी ट्यूरिंग मशीन~$G$ बना सकते हैं, जो एक ही आगत~$x$ लेती है, $H$ को
चलाकर तय करती है कि मशीन $M_x$, जिसका सूचकांक~$x$ है, आगत~$x$ पर रुकती है या
नहीं, और फिर इसके विपरीत करती है। दूसरे शब्दों में, यदि $H$ बताती है कि
$M_x$, आगत~$x$ पर रुकती है, तो $G$ अनंत चक्र में चली जाती है; और यदि
$H$ बताती है कि $M_x$, आगत~$x$ पर नहीं रुकती, तो $G$ रुक जाती है। क्या
$G$ अपने ही सूचकांक को आगत बनाने पर रुकती है? ऊपर का तर्क दिखाता है कि
वह तभी रुकती है जब वह नहीं रुकती---यह विरोधाभास है। अतः ऐसी ट्यूरिंग
मशीन~$H$ के अस्तित्व की हमारी मान्यता असत्य है।
\end{explain}
\end{tagblock}

\end{document}
